\documentclass{article}
\usepackage[utf8]{inputenc}
\usepackage[T1]{fontenc}
\usepackage{geometry}
\geometry{
    a4paper,
    left=2.5cm,
    right=2.5cm,
    top=2.5cm,
    bottom=2.5cm
}
\usepackage{graphicx}
\usepackage{float}
\usepackage{hyperref}
\usepackage[english]{babel}
\usepackage{amsmath}
\usepackage{amssymb}
\usepackage{array}
\usepackage{booktabs}
\usepackage{caption}
\usepackage{subcaption}
\usepackage{enumitem}
\usepackage{fancyhdr}
\usepackage{titlesec}
\usepackage{url}
\usepackage{textcomp}
\usepackage{xcolor}
\usepackage{listings}
\usepackage{multirow}
\usepackage{longtable}
\usepackage{tabularx}
\usepackage{times}

% Set up page style
\pagestyle{fancy}
\fancyhf{}
\rfoot{\thepage}
\renewcommand{\headrulewidth}{0pt}

% Title formatting
\titleformat{\chapter}[display]{\normalfont\bfseries}{\chaptertitle\thechapter}{0pt}{\Huge}
\titlespacing*{\chapter}{0pt}{50pt}{40pt}

% Document info
\title{Template Document}
\author{Academic Support Agent}
\date{\today}

\begin{document}
\maketitle

\begin{titlepage}
\thispagestyle{empty}

\begin{tikzpicture}[remember picture,overlay]
  \coordinate [above=0cm] (toppoint) at (current page.north);
  % Merki H sk la  slands.
  \node[anchor=north] (logonator) at ($(toppoint)+(0cm,-4.2cm)$) {
    \includegraphics[width=4.2cm]{./uilogo}
    };

  % Titill ritger ar
  \node[anchor=north, align=center] (titlenator) at ($(logonator.south)+(0cm,-3.0cm)$) {
    \huge \sffamily \bfseries \thesistitle
    };

  \coordinate [above=0cm] (bottompoint) at (current page.south);

  % Sta setning texta um deild og verkefnager 
  \node[anchor=base, 
  		minimum width=\paperwidth, 
        align=center] at ($(bottompoint)+(0cm,3.4cm)$)  
    {
    \parbox{
      1\paperwidth
      }{
      \centering
      \bfseries\sffamily \Large 
      \thesissubject \\
      \thesisfaculty \\
	  \thesisschool \\
	  H sk li  slands \\
	  \thesismonth~\thesisyear
      }
    };

  % Nafn h fundar
  \node[anchor=north, align=center] (authornator) at ($(titlenator.south)+(0cm,-1.0cm)$) {
    \normalfont \Large \sffamily \thesisauthor
	};
    
  % N fn lei beinenda
  \node[anchor=north, align=center] (advisornator) at ($(authornator.south)+(0cm,-1.0cm)$) {
	\parbox{
      1\paperwidth
      }{
      \centering
      \sffamily \Large Lei beinandi\\
      \thesistutors
      }
    };

  % Nafn h fundar
  \node[anchor=north, align=center] (authornator) at ($(advisornator.south)+(0cm,-1.0cm)$) {
    \normalfont \sffamily  \thesiscredits~ECTS verkefni fyrir \reportcourse
	};
    
\end{tikzpicture}

\newpage 

\end{titlepage}
\pagenumbering{roman}

\setcounter{page}{5}
\section*{\huge  tdr ttur}
H r kemur  tdr ttur    slensku sem er a  h marki 250 or . Reyni  a  koma  tdr ttum   eina bla s u en ef tv r bla s ur eru nau synlegar   seinni bla s a  tdr ttar a  hefjast   oddat lus u (h gri s u).
\vfill

\section*{\huge Abstract}
 tdr ttur   ensku sem er a  h marki 250 or .
\vfill \vspace*{1cm}

\newpage





\chapter*{Skammstafanir}
\addcontentsline{toc}{chapter}{Skammstafanir}
   essum kafla mega koma fram listar yfir skammstafanir og/e a breytuheiti. Gefi  kaflanum nafn vi  h fi, t.d. Skammstafanir e a Breytuheiti.  essum kafla m  sleppa ef hans er ekki  rf. \\

The section could be titled: Glossary, Variable Names, etc.

\chapter*{ akkir}
\addcontentsline{toc}{chapter}{ akkir}
   essum kafla koma fram  akkir til  eirra sem hafa styrkt ranns knina me  fj rframl gum, a st u e a vinnu. T.d. styrktarsj ir, fyrirt ki, lei beinendur, og a rir a ilar sem hafa   einhvern h tt a sto a  vi  ger  verkefnisins,  .m.t. vinir og fj lskylda ef vi   .  akkir byrja   oddat lus u (h gri s u).


\chapter{Inngangur}
\pagenumbering{arabic}
\setcounter{page}{1}
Ingangur skal vera  tarlegur og gefa yfirs n yfir  a  ranns knasvi  sem verkefni  n r til.  ar skal greint fr  tilgangi og bakgrunni verkefnisins, sk rt fr  st u  ekkingar   vi komandi ranns knasvi i, fyrri ranns knum og samhengi verkefnisins vi   r.   skal ger  grein fyrir  eim spurningum sem leitast er vi  a  svara me  ranns kninni.

\section{Almennar uppl singar}
Fyrirs gn 1 er kaflaheiti. Feitletri  fyrirs gn 1   20 pt Verdana. Hafi  54 pt loftun yfir og byrji  n ja oddat lus u (h gri s u). Hafi  12 pt loftun undir   undan texta og 24 pt loftun undir (samtals) ef beint   undan fyrirs gn 2.

  sta  Verdana m  velja samb rilegt steinskriftar (sans serif) letur   allar fyrirsagnir en allar fyrirsagnir skulu   vera rita ar me  s mu leturger .

Meginm lstexti skal skrifa ur   Times New Roman, me  leturst r  12 og einf ldu l nubili. M lsgreinar skulu lofta ar me  0 pt bili a  ofan og 12 pt bili a  ne an,  ar sem fyrirsagnir stilla loftun fyrir ne an sig (og  ar me  fyrir ofan texta).

Leita skal a  XX sem hluta af or i,  ar sem  a  merkir atri i sem h fundur  arf a  breyta.

  sta  Times New Roman m  velja samb rilegt prentletur (serif). Allt meginm l skal   rita  me  s mu leturger . 

Allur texti ritger ar skal rita ur me  einum lit, sv rtum. Undantekningar eru leyf ar innan mynda. Ekki nota  hyperlinks    texta, sem   ver ur bl r og/e a me  undirl nu.

Notast skal vi  2,5 cm sp ss u fyrir ofan og   ytri hli  (ekki kjalmegin). Vi  kj lin skal b ta 0,5 cm (Gutter) til a  hafa samtals 3,0 cm sp ss u. Ne st   bla s u skulu vera 1,5 cm fr  ne ri br n   bla s utal og sp ss an skal vera 3,0 cm fr  ne ri br n bla s u a  texta.

Bla s ur meginm ls byrja a  n merast   1   fyrstu bla s u fyrsta kafla me  arab skum st l. Bla s utali  er still vi  ytri br n og ne st   bla s u. Heimildir og vi aukar n merast einnig me  sama h tti. Fj ldi bla s na   ritger  skr ist sem bla s un mer  ftustu prentu u s u.

Byrja skal fyrirs gn 1 efst   h gri s u. H r m  hugsa s r a  setja stuttan texta sem inngang a  kaflanum  ur en fyrsti undirkafli byrjar.  a  getur hj lpa  lesanda a   tta sig   inntaki kaflans.

Nota m  sk letur   h fi til a  draga athygli a  texta. Noti  feitletur enn sjaldnar. Ekki nota undirstrika an texta   ritger inni.

Fari  sparlega   notkun footnote.  r skulu vera n mera ar og birtast ne st    eirri s u sem fyrst vitnar    r e a flj tlega  ar   eftir.

N meri  og v si    form lur eftir venjum fagsvi s.

Velja m  inndr tt fyrstu l nu m lsgreinar um 1 cm   sta  12 pt loftunar   milli m lsgreina, en    arf a  b ta 12 pt vi  loftun yfir fyrirs gnum sem gera r  fyrir a  12 pt loftun komi fr  lokum m lsgreinar. Ekki skal nota b i inndr tt og loftun.

\section{Fyrirs gn 2}
Fyrirs gn 2 er undirfyrirs gn. Feitletri  fyrirs gn 2   16 pt Verdana. Hafi  samtals 24 pt loftun yfir fyrirs gn 2. Noti  12 pt loftun fyrir ne an fyrirs gn 2.

\subsection{Fyrirs gn 3}
Fyrirs gn 3 er s asta n mera a undirfyrirs gnin. Feitletri  fyrirs gn 3   12 pt Verdana me  18 pt loftun yfir samtals  egar texti er fyrir ofan e a fyrirs gn 2. 
\subsubsection{Fyrirs gn 4}
12 pt Verdana, engin kaflan mer, birtist ekki   efnisyfirliti, 12 pt loftun yfir og 6 pt undir

Fyrirs gn 4 skal ekki n mera. Fyrirs gn 4 er ritu    venjulegu 12 pt Verdana og hefur minnst 12 pt loftun yfir ef h n er undir texta, en meiri loftun undir h rri fyrirs gnum sem st rist af ne ri loftunum  eirra fyrirsagna. Noti  6 pt loftun undir fyrirs gn 4.
Ekki nota fyrirsagnir   l gra  tl nu stigi en fyrirs gn 4.

\section{Um fors u, kj l og baks u}
Ekki skal nota b merki/log  fyrirt kja, samstarfsa ila e a styrktara ila   fors u/baks u e a annars sta ar   ritger . Ekki skal setja mynd   fors u ritger ar. 

  texta er hins vegar skylt og r tt a  geta samstarfsa ila og styrktara ila,  a  skal gert   kaflanum  akkir (Acknowledgments) e a   form la.

  baks u er heimilt a  setja nafn prentsmi ju,   mi ja  h gri-vinstri   bla s u og mi ja  upp-ni ur innan litabor ans. Nafn prentsmi ju skal rita  me  st rst 10 pt Verdana   venjulegu hv tu letri.

Merki H sk la  slands   fors u skal hafa 4,2 cm  verm l.  a  skal sta setja merki  4,2 cm fr  efri br n. Fr  ne ri br n merkis skulu minnst vera 3,0 cm   titil ritger ar.

Loftun fyrir ne an nafn h fundar skal vera 1 cm ni ur a  efstu br n litabor ans.

Megin hluti litabor ans ne st   fors u skal spanna 7,7 cm fr  ne ri br n bla s u   A4-formi, en   spannar hann lengra bil  ar sem t knmynd a albyggingar kemur fyrir.

Litabor inn skal n  yfir kj linn og baks una, og  ar spanna 7,7 cm fr  ne ri br n bla s u.

Litabor i BS ritger a er gr r litur H sk la  slands. Litak ar litarins eru: CMYK: 0 : 0 : 0 : 70; Pantone: Cool Gray 11 C; RGB: 90 : 91 : 94.

Litabor i meistararitger a er litur Verkfr i- og n tt ruv sindasvi s H sk la  slands, appels nugulur. Litak ar litarins eru: CMYK:  0 : 75 : 100 : 0; Pantone: 158 C; RGB: 236 : 78 : 34.

Litabor i doktorsritger a er litur H sk la  slands, d kk bl r. Litak ar litarins eru: CMYK: 100 : 57 : 0 : 40; Pantone: 295 C; RGB: 0 : 46 : 85.

  kj l ritger ar skal rita nafn h fundar, stuttan titil ritger ar (mest 50 sl g me  bilum) og  rtal ritger ar   einni l nu me  Verdana letri. Nafn h fundar og stuttur titill ritger ar koma   hv ta fl tinn, en  rtali  kemur   hv tu letri   litabor ann.

\section{Almennt um ritger asm }
G i ritger ar endurspegla ekki einungis g i ranns knarinnar (hermana, l kana, greininga, o.fl.), heldur einnig g i ritger asm ar.   s ara samhenginu skipta margir  ttir m li, t.d. uppbygging og s gufl i, framsetning   hugmyndum og ni urst um, m lfar og heimildavinna.  v  er mikilv gt a  nemendur kynni s r hvernig best s  a  standa a  undirb ningi og skrifum ritger ar. Nemendur ver a a  tileinka s r fagmannleg vinnubr g    heimildaskr ningu og tilv sunum   samr i vi  lei beinanda. H r eru d mi um tv r b kur sem nemendur geta stu st vi :
\begin{itemize}
 \item Fri r k H. J nsson og Sigur ur J. Gr tarsson (2007). \textit{Gagnfr ikver fyrir h sk lanema}, H sk la tg fan, Reykjav k.
 \item Ingibj rg Axelsd ttir og  runn Bl ndal (2006).  \textit{Handb k um ritun og fr gang},  M l og menning, Reykjav k.
\end{itemize}
Prufu heimildarv sun~\cite{knuth-fa}.
\section{Doktorsritger : Ritger  e a greinasafn}
Doktorsritger  getur veri  hvort sem er   formi ritger ar e a safns greina sem hafa birst e a  er  tla  a  birtist   ritr ndum ritum. 

Ef doktorsritger in er greinasafn skal  vallt fylgja inngangur sem setur fram heildst  markmi ,  tarlegt yfirlit og samantekt   verkinu. M lt er me  a    lok ritger ar,   undan vi aukum, fylgi samantekinn listi yfir heimildir verksins.

Ef einhver greinanna   ritger inni hafa me h funda skal framlag doktorsnemans  tsk rt s rstaklega   inngangi, s r   lagi krefst  a   tsk ringa ef doktorsneminn er ekki a alh fundur greinarinnar. 

R lagt er a  greinar  r sem birtast   safninu s u settar    eim st l sem h r er l st fyrir doktorsritger ir, en s u ekki bein afrit birtra greina,  ar sem setning og  tlit sl kra greina getur falli  undir vernd h fundarr ttar  tg fua ila.  tg fua ilar leyfa almennt a  texti greina birtist b i   grein og doktorsritger .  a  er    byrg  nemanda a  kanna sl kt og framfylgja  eim sam ykktum sem nemandi hefur fallist   vi  birtingu greina.
 
\chapter{Efni og a fer ir}
H r skal greint fr  efnivi  ranns knar, efnum sem notu  voru og uppruna  eirra eftir  vi sem vi   ,   er ekki nau synlegt a  geta uppruna efna sem algeng eru   ranns knastofum   vi komandi grein. A fer um skal l st  annig a  lesandi geti tileinka  s r a fer ina og endurteki  m lingu e a ranns kn me  fulln gjandi h tti. S u a fer ir  egar  ekktar skal v sa  til heimilda. Ef a fer  er breytt fr   v  sem l st hefur veri  skal gera sk ra grein fyrir breytingunni. N rri a fer  skal l st  tarlega.   skal geti  t lvuforrita sem notu  eru, til d mis vi  t lfr ilega  rvinnslu. A fer afr i ranns knar skal a  jafna i vera   samr mi vi   a  sem venja er   vi komandi fr asvi i. Ef val   a fer um getur orka  tv m lis skal  a  r kstutt.

\section{Listar}
H r   eftir er d mi um upptalningarlista. Listinn m  vera  ttari,  .e. a  a eins s  0 pt bil   milli atri a, en   skal hafa 12 pt bil   undan fyrsta atri i listans.
\begin{itemize}
 \item N mer 1
 \item N mer 2
 \item N mer 3
\end{itemize}
Ef fyrsta l na eftir upptalningu er framhald s mu m lsgreinar og fyrir ofan skal ekki setja inn vi b tarloftun (e a ef inndr ttur er nota ur, skal ekki nota inndr tt   sl kri l nu). 

\chapter{Myndir og t flur}
 essi kafli s nir d mi um notkun mynda, taflna og v sun    r.
\section{Myndir}
Myndatexti skal sta setja undir myndum og skrifast me  sk letri.
Setja skal au a aukal nu fyrir ofan myndir.
\begin{figure}[!htb]
\centering
\includegraphics[width=0.47\textwidth]{uilogo}
\caption[D mi um myndatexta (fyrir ne an mynd).]{D mi um myndatexta (fyrir ne an mynd)} \label{fig:Array}
\end{figure}
Mikilv gt er a  skilgreina myndir me  ,,paragraph format :  keep with next  til a  rj fa ekki tengsl   milli myndar og myndtexta. Mynd m  vera mi ju  og skal   einnig mi ja myndartextann. Letur innan myndar skal vera   steinskrift (sans serif), t.d. Verdana, og ekki minna en 10 pt. Tryggi  a  letur, t kn og l nur sj ist sk rt eftir  tprentun.

H gt er a  l ta n mera og merkja myndir sj lfvirkt me   v  a  gera Insert -- Reference -- Caption -- Mynd e a Tafla. Varist a  velja hyperlink. 

V sa m    mynd me   v  a  velja Insert -- Reference -- Cross-Reference -- Mynd e a Tafla. Varist a  velja hyperlink og velji  a  eins Label og Number. T.d. sj   essa tilv sun   Mynd 3.1 sem d mi. 

\section{T flur}
Einnig m  n mera t flur sj lfkrafa svipa  og myndir. Nota skal sk letur   t flutitil. Textinn skal standa fyrir ofan t flu og fylgja t flunni.  Ekki nota tv falt l nubil e a hafa space before   t flum. Meginreglan vi  t fluger  er a  hafa  r einfaldar og eins f  strik og m gulegt er. Tafla m  vera mi ju    bla s u og skal   l ta t flutitil byrja vi  vinstri br n t flu.


\begin{table}[htb]
\centering
\caption{D mi um t flutexta (fyrir ofan t flu).}
     \sffamily \begin{tabularx}{1.0\textwidth}{ p{5cm}  p{5cm}  p{5cm} }
    \hline
   \textbf{Taflan} \hfill & \textbf{Er} \hfill & \textbf{Eins} \\ \hline
    Og & H n & g ti\\
    Liti  &  t &   ritger \\ \hline
    \end{tabularx} \normalfont
\label{table:Emissivity}
\end{table}

Almennt skal ekki nota loftun fyrir ne an texta   t flu, og stylla loftun fyrir ne an   0 pt.
Mikilv gt er a  skilgreina t flutexta me  ,,format paragraph: keep with next og keep lines together  til a  rj fa ekki tengsl   milli t flutexta og t flu. Ef tafla er mj g l ng m  klj fa hana   milli bla s na og   ver ur a  setja (\textit{Framhald})   aukal nu beint fyrir ne an t fluna, h gri stillt vi  h gri br n t flu.

D mi um sj lfvirka tilv sun   t flu, bara nota Label and Number, ekki nota hyperlink e a caption text. T.d. Tafla 3.1 s nir d mi um t flu.

\chapter{Ni urst ur}
A  jafna i skal h r einungis greina fr  ni urst um ranns kna e a m linga.   styttri ritger , t.d. lokaverkefni til BS-pr fs, getur komi  til greina a  sameina ni urst ur og umr ur   einn kafla.

Kappkosta  skal a  setja ni urst ur fram   skipulegan og sk ran h tt,  n m lalenginga. Einungis skal sk rt fr   v  sem telja m  sta reyndir. Ni urst ur skulu eftir  v  sem unnt er settar fram   formi taflna e a mynda og innihald  eirra sk rt og meginatri i dregin fram   texta. Myndir og t flur skulu settar inn   texta  ar sem um  r er r tt. Myndir skulu vera sk rar og st rar, a  jafna i skeytt inn   skjali  rafr nt. L nurit skulu fullfr gengin  annig a   sar s u sk rt merktir og t kn sk r . Myndatexti me  n meri og heiti myndar skal fylgja ne an vi  hverja mynd og skal hann vera stuttur en   l sa vi fangsefni myndar   fulln gjandi h tt. Sama gildir um t flur, sem skulu vera sk rar og greinarg ar. Texti me  n meri og heiti t flu skal fylgja.

\section{Efni s tt fr   ru skjali}
D mi um heimild~\cite{dirac} og l ka~\cite{einstein}.
% Content from YtraEfniFyrirInclude would be included here


\chapter{Umr ur/skil}
   essum kafla er ni urst um ger  skil og  r r ddar og sk r ar. Ger  er grein fyrir einst kum  ttum ranns knanna,  r sk r ar og settar   samhengi. Gera skal grein fyrir  g llum og  vissu ttum eftir  v  sem vi   . Ef vi    eru ni urst ur bornar saman vi  ni urst ur annarra ranns kna og settar fram hugmyndir e a till gur um frekari ranns knir.

\chapter{ lyktanir/lokaor }
H r eru settar fram   stuttu m li helstu  lyktanir sem draga m  af ni urst um ranns knarinnar.


%
% Heimildir
%
\printbibliography[heading=bibintoc] % B tir vi  heimildum og f rir   efnisyfirlit

A alfyrirs gn heimildaskr r skal birtast   efnisyfirliti, h n skal hafa sama form og fyrirs gn eitt en vera  n kaflan mers. Eins og a rar a alfyrirsagnir skulu heimildir byrja   n rri bla s u og h n skal vera oddat lu (h gri) s a.
Noti  eitt samr mt form   heimildaskr  og tilvitnunum   ritger inni. Velji  form   heimildaskr  og tilvitnunum   samr i vi  lei beinanda til  ess a  venjur fagsvi s ver i uppfylltar.


%
% Vi hengi
%

\renewcommand{\chaptername}{Vi auki}
\chapter{Vi auki}
Gagnlegir punktar og  bendingar
\begin{itemize}
 \item BS ritger ir skulu prenta ar   A4.
 \item Meistararitger ir skulu prenta ar   A4.
 \item  tg fust r  doktorsritger a er B5.
\begin{itemize}
 \item[-] I ulega er doktorsritger  unnin   A4 en svo sm kku  vi  prentun. Mikilv gt er a  tryggja a  allur texti, m.a.   myndum og t flum, sj ist sk rt eftir sm kkun   B5.  v  ver ur a  sko a dr g af ritger  eftir sm kkun   B5  ur en ritger in er prentu    fj lda eintaka.
\end{itemize}

\end{itemize}

\begin{itemize}
 \item Mi a skal vi  a  ritger  s  prentu  b u megin   bla s ur og byrja skal alla kafla   h gri s u opnu.
 \item Hafa skal samband vi  prentsmi ju  ur en handriti er skila . Oft vilja  essir a ilar f  ritger ina   PDF-formi. Hafa skal   huga a  litprentun er mun d rari en svart/hv t prentun. 
 \item Til eru mismunandi g i/upplausn   PDF-skj lum. Prentsmi jur bi ja gjarnan um h g a upplausn / prentunarupplausn sem er meiri en ,,venjuleg  PDF-upplausn sem notu  er vi  skj l sem vistu  eru   netinu.  etta er stillingaratri i  ur en PDF-skjal er b i  til. Almennt vilja prentsmi jur h stu m gulegu upplausn. Einnig vilja prentsmi jur almennt a  allt letur og myndir s u skilgreindar ( Embedded ) innan   PDF skjalinu.
\end{itemize}
\end{document}